\subsection{Interakce učitele s návrhy na řádcích}
\label{subsec:implementation-interaction}

U každého jednotlivého navrhovaného komentáře má učitel k dispozici čtyři akce.

\begin{itemize}
    \item \textbf{Přijmout} (\textit{accept}) znamená, že návrh je schválen a systém automaticky vytvoří standardní komentář viditelný pro studenta.
    Přijatý návrh zůstává v databázi propojen s vytvořeným komentářem, což umožňuje zpětné dohledání vztahu mezi LLM návrhem a finálně publikovaným komentářem.

    \item \textbf{Upravit} (\textit{edit}) umožňuje učiteli pozměnit obsah návrhu před jeho publikací.
    Komentář se poté vytvoří stejným způsobem jako při přijetí, avšak s upraveným textem.

    \item \textbf{Zamítnout} (\textit{reject}) označí návrh jako nevhodný, takže se studentovi nezobrazí.
    Zamítnutý návrh zůstává v databázi pro pozdější analýzu kvality modelu, avšak žádný odpovídající standardní komentář není vytvořen.

    \item \textbf{Ohodnotit} (\textit{rate}) umožňuje učiteli poskytnout zpětnou vazbu o kvalitě a relevanci návrhu, přičemž hodnocení lze zadat i při přijetí nebo zamítnutí.
    Podrobněji je systém hodnocení popsán níže.
\end{itemize}

Před publikací návrhu je vhodné, aby učitel zvážil jeho obsah, protože navrhované komentáře mohou obsahovat nepřesnosti nebo být v daném kontextu zcela nevhodné.
Hodnocení podle metrik definovaných v \secref[sekci]{subsec:criteria-quality} tvoří základ zpětné vazby, která může v budoucnu sloužit k iterativnímu zlepšování promptů.
Panel pro přijetí nebo zamítnutí návrhu včetně hodnocení kvality a relevance hvězdičkami je viditelný v pravém horním rohu komentáře na \figref[obrázku]{fig:frontend-comment}.

\endinput